\documentclass[12pt,oneside]{book}

% ==========================================
% METADATOS Y VARIABLES
% ==========================================
\newcommand{\universidad}{Universidad de Buenos Aires (UBA)}
\newcommand{\facultad}{Facultad de Derecho}
\newcommand{\carrera}{Abogacía}
\newcommand{\materia}{Derecho Procesal Civil}
\newcommand{\titulo}{Nociones Fundamentales del Derecho Procesal Civil}
\newcommand{\autor}{Isaac Edgar Camacho Ocampo}
\newcommand{\anio}{2026}
\newcommand{\reflexion}{Comprender los conceptos es el primer paso para convertir el estudio en conocimiento.}

% ==========================================
% PAQUETES OBLIGATORIOS Y CONFIGURACIÓN
% ==========================================
\usepackage[utf8]{inputenc}
\usepackage[T1]{fontenc}
\usepackage[spanish,es-nodecimaldot]{babel}

\usepackage[
    top=2.5cm,
    bottom=2.5cm,
    left=3cm,
    right=2.5cm,
    headheight=15pt
]{geometry}

\usepackage{amsmath}
\usepackage{amssymb}
\usepackage{mathtools}

\usepackage{graphicx}
\usepackage{float}
\usepackage{caption}
\usepackage{subcaption}

\usepackage{booktabs}
\usepackage{array}
\usepackage{longtable}
\usepackage{tabularx}

\usepackage{enumitem}
\usepackage{xcolor}
\usepackage[most]{tcolorbox}
\usepackage{fancyhdr}
\usepackage{ragged2e}

\usepackage[
    colorlinks=true,
    linkcolor=blue!50!black,
    citecolor=green!40!black,
    urlcolor=blue!60!black,
    pdftitle={\titulo},
    pdfauthor={\autor},
    pdfsubject={Apuntes universitarios},
    bookmarks=true
]{hyperref}

% ==========================================
% RUTAS DE IMÁGENES
% ==========================================
\graphicspath{
    {./img/}
    {./graficos/}
    {./figuras/}
}

% ==========================================
% CONFIGURACIÓN GENERAL
% ==========================================
\setcounter{tocdepth}{2}
\setlength{\parindent}{0pt}
\setlength{\parskip}{0.5em}

\setlist[itemize]{
    topsep=0.3em,
    itemsep=0.2em
}

\setlist[enumerate]{
    topsep=0.3em,
    itemsep=0.2em
}

\clubpenalty=10000
\widowpenalty=10000

% ==========================================
% ENCABEZADOS Y PIES DE PÁGINA
% ==========================================
\pagestyle{fancy}
\fancyhf{}
\fancyhead[L]{\nouppercase{\leftmark}}
\fancyhead[R]{\materia}
\fancyfoot[C]{\thepage}

\renewcommand{\headrulewidth}{0.4pt}
\renewcommand{\footrulewidth}{0pt}

\fancypagestyle{plain}{
    \fancyhf{}
    \fancyfoot[C]{\thepage}
    \renewcommand{\headrulewidth}{0pt}
}

% ==========================================
% COMANDOS MATEMÁTICOS Y AUXILIARES
% ==========================================
\newcommand{\abs}[1]{\left|#1\right|}
\newcommand{\norm}[1]{\left\lVert#1\right\rVert}
\newcommand{\vect}[1]{\mathbf{#1}}
\newcommand{\deriv}[2]{\frac{d#1}{d#2}}
\newcommand{\pderiv}[2]{\frac{\partial#1}{\partial#2}}

% ==========================================
% CAJAS ACADÉMICAS (TCOLORBOX)
% ==========================================
\newtcolorbox{definition}[1]{
    enhanced,
    breakable,
    title={Definición: #1},
    fonttitle=\bfseries,
    colback=blue!3,
    colframe=blue!50!black,
    boxrule=0.8pt,
    arc=2mm
}

\newtcolorbox{important}{
    enhanced,
    breakable,
    title={Importante},
    fonttitle=\bfseries,
    colback=yellow!8,
    colframe=orange!70!black,
    boxrule=0.8pt,
    arc=2mm
}

\newtcolorbox{example}[1]{
    enhanced,
    breakable,
    title={Ejemplo: #1},
    fonttitle=\bfseries,
    colback=green!4,
    colframe=green!40!black,
    boxrule=0.8pt,
    arc=2mm
}

\newtcolorbox{formula}[1]{
    enhanced,
    breakable,
    title={Fórmula / Regla: #1},
    fonttitle=\bfseries,
    colback=purple!4,
    colframe=purple!50!black,
    boxrule=0.8pt,
    arc=2mm
}

\newtcolorbox{exercise}[1]{
    enhanced,
    breakable,
    title={Ejercicio: #1},
    fonttitle=\bfseries,
    colback=gray!5,
    colframe=gray!60!black,
    boxrule=0.8pt,
    arc=2mm
}

\newtcolorbox{note}{
    enhanced,
    breakable,
    title={Nota},
    fonttitle=\bfseries,
    colback=white,
    colframe=gray!50,
    boxrule=0.6pt,
    arc=2mm
}

% ==========================================
% INICIO DEL DOCUMENTO
% ==========================================
\begin{document}

% ------------------------------------------
% PORTADA
% ------------------------------------------
\begin{titlepage}
\thispagestyle{empty}
\begin{center}

\vspace*{1cm}

{\large \textsc{\universidad}}

\vspace{0.2cm}

{\large \textsc{\facultad}}

\vspace{0.2cm}

{\large \carrera}

\vspace{3cm}

{\Huge \textbf{\materia}}

\vspace{1cm}

{\LARGE \titulo}

\vspace{0.8cm}

\textit{\reflexion}

\vspace{1.2cm}

\rule{0.70\textwidth}{0.5pt}

\vspace{0.5cm}

{\large Apuntes de estudio}

\vspace{0.5cm}

\rule{0.70\textwidth}{0.5pt}

\vfill

{\large \autor}

\vspace{0.3cm}

{\large \anio}

\end{center}

\vfill

\begin{flushleft}
\textit{Este es un modesto aporte a los estudiantes de ``\materia'' no pretende sustituir el material oficial de la materia.}
\end{flushleft}
\end{titlepage}

% ------------------------------------------
% ÍNDICE
% ------------------------------------------
\frontmatter
\tableofcontents
\mainmatter

% ------------------------------------------
% CONTENIDO ACADÉMICO
% ------------------------------------------
\chapter{Introducción y Fundamentos del Derecho Procesal Civil}

\section{Presentación de la Cátedra y Dinámica de Estudio}
El estudio del Derecho Procesal Civil requiere comprender el funcionamiento dinámico del sistema judicial más allá de la mera memorización normativa. Resulta fundamental abordar la materia mediante un análisis conceptual de las instituciones procesales, complementado con la práctica forense en la elaboración de demandas, contestaciones e interrogatorios de testigos.

Para la preparación de la asignatura se toma como referencia la obra \textit{Nociones Básicas del Derecho Procesal Civil}, así como los tratados y manuales de autores clásicos de la doctrina (Palacio, Falcón, Arazi) y el cuerpo normativo dado por el Código Procesal Civil y Comercial de la Nación acompañado de sus leyes complementarias.

\section{El Conflicto Jurídico y el Origen del Proceso}
Para comprender la razón de ser del Derecho Procesal, es preciso trazar una distinción clara entre el derecho sustancial (o de fondo) y el derecho procesal:
\begin{itemize}
    \item \textbf{Derecho Sustancial:} Regula las relaciones jurídicas ordinarias entre los sujetos dentro de la sociedad (por ejemplo, los derechos y obligaciones derivados del Código Civil y Comercial de la Nación).
    \item \textbf{Derecho Procesal:} Entra en funcionamiento cuando la relación jurídica sustancial entra en conflicto, las normas de fondo dejan de ser operativas de manera voluntaria entre las partes y se requiere la intervención del Estado para resolver la controversia.
\end{itemize}

\begin{example}{Agotamiento de la relación sustancial y surgimiento del proceso}
Una persona aborda un transporte público celebrando un contrato de transporte (regido por el derecho sustancial). A raíz de un accidente, el pasajero sufre lesiones e internación. Al quedar insatisfecha la pretensión de reparación y agotada la vía directa entre las partes, surge una nueva relación jurídica de carácter procesal en la que interviene un tercero imparcial: el juez, bajo la regulación del Código Procesal Civil y Comercial.
\end{example}

\begin{example}{Conflicto en relaciones de consumo e infortunios laborales}
Un consumidor adquiere un bien defectuoso y la garantía no responde; ante la falta de solución directa, el conflicto debe encauzarse mediante la vía procesal ante la justicia ordinaria o los tribunales especializados de consumo. Del mismo modo, ante un accidente laboral derivado de fallas de seguridad, la relación laboral de fondo entra en crisis y la controversia se traslada al ámbito procesal correspondiente.
\end{example}

\begin{definition}{Sentencia judicial}
La \textbf{sentencia} es la norma individual dictada por el órgano jurisdiccional para resolver el conflicto entre las partes. Funciona como una norma creada a medida del caso concreto que restituye la operatividad del derecho sustancial por mandato judicial.
\end{definition}

\section{Autonomía e Unidad del Derecho Procesal}
El Derecho Procesal es una disciplina científica autónoma e independiente, dotada de codificación propia, doctrina especializada y principios estructurantes. Sin perjuicio de las diversas especialidades, se sostiene la \textbf{unidad del Derecho Procesal}, en tanto todas las ramas comparten un mismo pedestal conceptual integrado por la \textbf{jurisdicción}, la \textbf{acción/pretensión} y el \textbf{proceso}.

Entre las ramas especializadas se destacan:
\begin{itemize}
    \item \textbf{Derecho Procesal Penal:} Articula el ejercicio de la acción penal pública mediante el Ministerio Público, el régimen de medidas cautelares (como la prisión preventiva) y la vía recursiva específica (casación, nulidad, inconstitucionalidad).
    \item \textbf{Derecho Procesal Laboral:} Se estructura a partir del principio protectorio, destinado a compensar la asimetría entre trabajador y empleador en el desarrollo del proceso.
    \item \textbf{Derecho Procesal Constitucional:} Aborda las garantías de tutela directa como el amparo, la acción declarativa de inconstitucionalidad y el recurso extraordinario federal.
    \item \textbf{Derecho Procesal Administrativo:} Regula el control jurisdiccional sobre los actos del Poder Ejecutivo y la administración pública.
\end{itemize}

\section{Fuentes del Derecho Procesal}
Las normas e instituciones del Derecho Procesal se alimentan de las siguientes fuentes fundamentales:

\begin{enumerate}
    \item \textbf{Constitución Nacional:} Constituye la fuente primaria. El artículo 18 CN consagra el principio del juez natural y el debido proceso, mientras que la parte orgánica establece la estructura del Poder Judicial.
    \item \textbf{Tratados Internacionales de Derechos Humanos:} Gozan de jerarquía constitucional en los términos del art. 75 inc. 22 de la Constitución Nacional.
    \begin{important}
    En el fallo \textit{Ekmekdjian c/ Sofovich}, la Corte Suprema de Justicia de la Nación reconoció la operatividad directa de las cláusulas de los tratados internacionales de derechos humanos, estableciendo que no requieren de una ley reglamentaria interna para ser exigibles ante los tribunales.
    \end{important}
    \item \textbf{Leyes Nacionales y Provinciales:} Códigos de procedimiento y leyes orgánicas del Poder Judicial.
    \item \textbf{Decretos de Necesidad y Urgencia (DNU):} Dictados por el Poder Ejecutivo y sometidos al control parlamentario de la Comisión Bicameral (Ley 26.122). Su incidencia en el proceso se evidencia cuando disponen la suspensión de plazos, ejecuciones o lanzamientos.
    \item \textbf{Jurisprudencia:} Emanada principalmente de los fallos de la Corte Suprema de Justicia de la Nación y los tribunales de alzada.
    \item \textbf{Doctrina:} Estudios y construcciones teóricas de los autores especializados que interpretan y sistematizan la norma procesal.
\end{enumerate}

\section{Evolución Histórica de los Sistemas Procesales}
El estudio de las instituciones procesales exige comprender su trayectoria histórica y la confluencia de diferentes tradiciones jurídicas:

\begin{itemize}
    \item \textbf{Proceso Romano:} Caracterizado por el desarrollo en dos fases: la etapa \textit{in iure} (ante el magistrado) y la etapa \textit{apud iudicem} (ante el juez o árbitro).
    \item \textbf{Derecho Germánico y Edad Media:} Invasiones bárbaras que introdujeron elementos místico-religiosos como las ordalías o "juicios de Dios". Posteriormente, el influjo del Derecho Canónico consolidó el denominado "derecho común", el cual llegó al ámbito hispanoamericano a través de las Leyes de Enjuiciamiento Civil españolas de 1855 y 1881.
    \item \textbf{Sistemas de Enjuiciamiento Comparados:}
    \begin{itemize}
        \item \textbf{Sistema Continental (\textit{Civil Law}):} Modelo adoptado por la Argentina, con marcada tradición escrita, centralidad de la prueba documental y ordenación del proceso mediante el expediente.
        \item \textbf{Sistema Anglosajón (\textit{Common Law}):} Centrado en la oralidad, el juicio por jurados y el predominio de la prueba testimonial.
    \end{itemize}
\end{itemize}

El Código Procesal Civil y Comercial de la Nación, sancionado en 1967 bajo la influencia de la doctrina procesal italiana, constituye la base del ordenamiento procesal nacional actual.

\chapter{El Pedestal Básico del Derecho Procesal}

\section{La Jurisdicción}

\begin{definition}{Jurisdicción}
Proveniente del latín \textit{iuris dictio} ("decir el derecho"), la \textbf{jurisdicción} es la potestad y función pública emanada de la Constitución Nacional de la que están investidos los órganos del Poder Judicial para resolver conflictos mediante el dictado de una norma individual (sentencia) con fuerza de cosa juzgada.
\end{definition}

\subsection{Distinción entre Jurisdicción y Competencia}
Resulta indispensable no confundir ambos conceptos:
\begin{itemize}
    \item \textbf{Jurisdicción:} Función constitucional unitaria e indivisible de impartir justicia.
    \item \textbf{Competencia:} Capacidad o aptitud otorgada a cada juez para ejercercer la jurisdicción en un caso determinado. Es la distribución práctica del trabajo judicial establecida por el legislador según diversos criterios:
    \begin{itemize}
        \item \textit{Materia:} Civil, comercial, penal, laboral, contencioso administrativo.
        \item \textit{Territorio:} Delimitación geográfica de la sede judicial.
        \item \textit{Grado:} Instancias judiciales (juzgados de primera instancia, cámaras de apelación, CSJN).
        \item \textit{Turno:} Asignación temporal de causas.
    \end{itemize}
\end{itemize}

\subsection{Los Cinco Momentos de la Jurisdicción}
El ejercicio de la función jurisdiccional despliega cinco facultades o momentos clave:

\begin{enumerate}
    \item \textbf{Notio:} Facultad de conocer y tomar cognición del conflicto presentado en la demanda.
    \item \textbf{Vocatio:} Potestad de convocar y citar a las partes para que comparezcan al proceso.
    \item \textbf{Coertio:} Poder de coerción para emplear los medios necesarios que garanticen el normal desarrollo del proceso (ej. ordenar el comparendo de un testigo por la fuerza pública).
    \item \textbf{Iudicium:} Facultad decisoria central; se concreta al dictar la sentencia definitiva.
    \item \textbf{Executio:} Potestad de ejecutar la sentencia y hacer cumplir lo resuelto mediante el uso de la fuerza pública si el condenado no lo realiza voluntariamente (ej. firma del juez en sustitución del deudor para escriturar un bien).
\end{enumerate}

\subsection{Imperio y Jurisdicción en Ámbitos No Judiciales}
\begin{important}
No debe equipararse la jurisdicción con el \textbf{imperio}. El \textit{imperio} es el poder de coerción y uso de la fuerza pública, atributo que comparten los tres poderes del Estado. La \textit{jurisdicción} (\textit{iuris dictio}), en cambio, es la potestad exclusiva de resolver controversias con autoridad de cosa juzgada, facultad privativa del Poder Judicial.
\end{important}

Existen manifestaciones jurisdiccionales singulares fuera del esquema judicial ordinario:
\begin{itemize}
    \item \textbf{Jurisdicción Administrativa:} Ejercida por organismos públicos (ej. AFIP, entes reguladores). La CSJN fijó en la doctrina de los fallos \textit{Fernández Arias c/ Poggio} y \textit{Ángel Estrada y Cía.} que las decisiones administrativas solo son válidas si están sujetas a un control judicial suficiente y no imponen condenas resarcitorias propias de la justicia ordinaria.
    \item \textbf{Jurisdicción Arbitral:} Jueces privados seleccionados voluntariamente por las partes. Poseen jurisdicción para dictar el laudo, pero carecen de \textit{imperio} para ejecutarlo, debiendo recurrir a la justicia ordinaria para su cumplimiento forzoso.
\end{itemize}

\section{Acción y Pretensión}

\subsection{La Polémica Windscheid--Muther (1856) y la Autonomía de la Acción}
El debate entre Bernhard Windscheid y Theodor Muther en 1856 constituye el hito histórico fundacional del Derecho Procesal como ciencia autónoma. Hasta entonces se concebía a la \textit{actio} romana como el mismo derecho sustancial "en pie de guerra".

La polémica demostró que el derecho de \textbf{acción} es abstracto y autónomo del derecho material subyacente: un sujeto tiene derecho a accionar y acudir ante la jurisdicción independientemente de que la sentencia final declare con o sin lugar su reclamo. Posteriormente, Oskar von Bülow (1868) y Giuseppe Chiovenda (1913) consolidaron el carácter científico del proceso como una relación jurídica autónoma.

\subsection{Delimitación entre Acción, Pretensión y Demanda}
\begin{itemize}
    \item \textbf{Acción:} Derecho constitucional genérico de peticionar ante las autoridades (art. 14 CN). Es una noción abstracta.
    \item \textbf{Pretensión:} La concreción de la acción orientada a un conflicto determinado. Es la declaración de voluntad formal dirigida al juez y frente al demandado para que se satisfaga un interés concreto.
    \item \textbf{Demanda:} El acto procesal instrumental que contiene en su texto la pretensión e inicia formalmente el proceso.
\end{itemize}

El planteo de la pretensión configura la relación jurídica procesal representada en el triángulo procesal:

\begin{formula}{El Triángulo Procesal}
\[
\text{Actor (Pretende)} \longrightarrow \text{Juez (Tercero Imparcial)} \longleftarrow \text{Demandado (Resiste)}
\]
La relación jurídica es bilateral y se constituye entre las dos partes litigantes y el órgano jurisdiccional.
\end{formula}

\section{El Proceso Judicial como Sistema}

\subsection{Concepción Sistémica del Proceso}
Conforme a la Teoría General de Sistemas (von Bertalanffy), el proceso no se agota en una suma de trámites, sino que constituye un verdadero \textbf{sistema} coordinado de elementos interconectados para alcanzar un objetivo común: la resolución del conflicto mediante una sentencia justa.

\begin{figure}[H]
    \centering
    \caption{Componentes del Sistema Procesal}
    \begin{itemize}
        \item \textbf{Insumos:} Sujetos procesales (actor, demandado, juez, auxiliares), recursos económicos, infraestructura e instrumentos normativos y tecnológicos.
        \item \textbf{Procesador:} Conjunto ordenado de etapas procedimentales.
        \item \textbf{Resultado:} La sentencia judicial (norma individual).
        \item \textbf{Mecanismo de Retroalimentación:} Los recursos procesales.
    \end{itemize}
\end{figure}

\subsection{Etapas del Procesador}
El desarrollo del proceso judicial transcurre a través de cuatro etapas concatenadas:
\begin{enumerate}
    \item \textbf{Etapa Introductoria:} Constitución de la litis mediante la presentación de la demanda, citación y contestación de la demanda.
    \item \textbf{Etapa Probatoria:} Ofrecimiento, admisión y producción de las medidas probatorias para acreditar los hechos controvertidos.
    \item \textbf{Etapa Conclusional:} Evaluación del material probatorio a través de los alegatos de bien probado.
    \item \textbf{Etapa Decisoria:} Emisión de la sentencia por parte del juez.
\end{enumerate}

\subsection{La Retroalimentación mediante el Sistema Recursivo}
El mecanismo de \textbf{retroalimentación} actúa como un control de calidad del resultado del sistema procesal. Ante la interposición de un recurso procesal (ej. apelación), el tribunal de alzada revisa la decisión del inferior. Si la sentencia es revocada o modificada, la corrección orienta las futuras decisiones del juzgador y la reformulación de estrategias de los litigantes.

\section{División de Poderes y Función Jurisdiccional}
El Poder del Estado es uno solo y emana del pueblo. Para su ejercicio se organiza en tres departamentos funcionales:
\begin{itemize}
    \item \textbf{Poder Legislativo:} Función atípica y primaria de sancionar la ley general.
    \item \textbf{Poder Ejecutivo:} Función de administración e implementación.
    \item \textbf{Poder Judicial:} Función esencial de dirimir conflictos aplicando la voluntad de la ley a casos concretos.
\end{itemize}

A pesar de esta especialización, existen interrelaciones y funciones cruzadas (ej. la emisión de Acordadas por la CSJN reglamentando la notificación electrónica posee naturaleza normativa interna). No obstante, la función jurisdiccional estricta —la potestad irrevocable de juzgar— recae privativamente en el Poder Judicial.

% ------------------------------------------
% QUADRO CORNELL DE REPASO
% ------------------------------------------
\chapter{Cuadro Cornell de Repaso General}

\section{Tabla de Repaso Integrada}

\begin{table}[H]
\centering
\begin{tabularx}{\textwidth}{
    >{\raggedright\arraybackslash}p{0.30\textwidth}
    >{\raggedright\arraybackslash}X
}
\toprule
\textbf{Preguntas / Palabras clave} & \textbf{Notas / Respuestas} \\
\midrule

\textbf{¿Qué es el Derecho Procesal Civil y cuál es su finalidad?} &
Es la disciplina que rige la actividad jurisdiccional y los trámites para resolver un conflicto sustancial insatisfecho, restituyendo la operatividad del derecho de fondo mediante la sentencia. \\

\textbf{¿Qué diferencia existe entre el derecho sustancial y el procesal?} &
El derecho sustancial regula la vida pacífica y cotidiana (ej. Código Civil y Comercial); el procesal se activa cuando surge el conflicto y se requiere la intervención del juez. \\

\textbf{¿Qué es una sentencia judicial?} &
Es la norma individual dictada por el juez que resuelve de forma imperativa el conflicto planteado entre las partes en un caso concreto. \\

\textbf{¿A qué se refiere la unidad del Derecho Procesal?} &
Implica que, a pesar de existir ramas especializadas (penal, laboral, administrativo), todas se fundamentan en las mismas instituciones básicas: jurisdicción, acción/pretensión y proceso. \\

\textbf{¿Cuáles son las fuentes principales del Derecho Procesal?} &
La Constitución Nacional (art. 18), Tratados Internacionales (art. 75 inc. 22), leyes procesales, DNU, jurisprudencia y doctrina. \\

\textbf{¿Qué significación tuvo el fallo \textit{Ekmekdjian c/ Sofovich}?} &
Estableció la operatividad directa de los Tratados Internacionales de Derechos Humanos sin necesidad de reglamentación legislativa previa. \\

\textbf{¿Qué es la jurisdicción y cómo se distingue de la competencia?} &
La jurisdicción es la potestad constitucional unitaria de impartir justicia. La competencia es el límite o porción de la jurisdicción asignada a cada juez por materia, territorio, grado o turno. \\

\textbf{¿Cuáles son los cinco momentos de la jurisdicción?} &
\textbf{1.~Notio} (conocer), \textbf{2.~Vocatio} (convocar), \textbf{3.~Coertio} (compeler), \textbf{4.~Iudicium} (juzgar) y \textbf{5.~Executio} (ejecutar la sentencia). \\

\textbf{¿En qué se diferencian el imperio y la jurisdicción?} &
El imperio es la facultad de usar la fuerza pública (presente en los tres poderes). La jurisdicción es la potestad exclusiva del Poder Judicial de resolver controversias con fuerza de cosa juzgada. \\

\textbf{¿Qué validez tiene la jurisdicción administrativa y arbitral?} &
La administrativa requiere siempre control judicial posterior (\textit{Fernández Arias}). La arbitral posee jurisdicción pero no imperio; para ejecutar el laudo debe acudir a los jueces estatales. \\

\textbf{¿Qué trascendencia tuvo la polémica Windscheid--Muther (1856)?} &
Consolidó la autonomía del Derecho Procesal al demostrar que el derecho de acción es independiente y abstracto respecto del derecho sustancial reclamado. \\

\textbf{¿Cómo se relacionan acción, pretensión y demanda?} &
La acción es el derecho abstracto de peticionar; la pretensión es la concreción de ese reclamo frente a un demandado; la demanda es el escrito que la formaliza. \\

\textbf{¿Cómo se concibe el proceso desde la Teoría de Sistemas?} &
Como un sistema dinámico compuesto por insumos (sujetos y recursos), procesador (etapas de la litis), resultado (sentencia) y retroalimentación (recursos procesales). \\

\textbf{¿Cuáles son las cuatro etapas del proceso?} &
1.~Etapa introductoria, 2.~Etapa probatoria, 3.~Etapa conclusional (alegatos) y 4.~Etapa decisoria (sentencia). \\

\midrule

\multicolumn{2}{p{\textwidth}}{
\textbf{Resumen del Capítulo:}
El Derecho Procesal Civil es una ciencia autónoma y unitaria que se fundamenta en la trilogía estructural compuesta por la \textbf{jurisdicción} (potestad del Estado de juzgar y hacer ejecutar lo juzgado desplegada en cinco momentos: \textit{notio}, \textit{vocatio}, \textit{coertio}, \textit{iudicium} y \textit{executio}), la \textbf{pretensión} (concreción procesal de la acción autónoma encausada mediante la demanda) y el \textbf{proceso} (entendido como un sistema articulado en etapas que concluye con la sentencia y se autocorrige mediante el régimen recursivo).
} \\

\bottomrule
\end{tabularx}
\end{table}

\end{document}